\documentclass[12pt,brazil,a4paper, twocolumn]{article}
\usepackage[T1]{fontenc}
\usepackage[utf8]{inputenc}
\usepackage{geometry}
\usepackage{graphicx}
\geometry{verbose,tmargin=2cm,bmargin=2cm,lmargin=2cm,rmargin=2cm}
\RequirePackage{times}
\usepackage{latexsym}
\usepackage{amsmath}
\usepackage{algorithm}
\usepackage{float}
\usepackage{algpseudocode}
\usepackage[hidelinks]{hyperref}
\usepackage{babel}

\begin{document}

\title{Interpolação de Newton -- Estimativa da altura máxima de um projétil}

\author{Agatha Schneider\\
Métodos Numéricos}

\date{\today}
\maketitle

\begin{abstract}
Este artigo tem como objetivo apresentar a resolução de um problema de interpolação numérica. O problema consiste em estimar a altura máxima atingida por um projétil a partir de uma tabela com alguns valores de tempo e altura. Para isso, foi utilizado o método de interpolação de Newton por diferenças divididas.

Além da resolução implementada, também foi feita uma segunda estimativa utilizando o ChatGPT. A ferramenta recebeu a mesma tabela de dados e foi solicitado que estimasse o tempo e a altura máxima do projétil. Ao final, os resultados foram comparados e discutidos.
\end{abstract}

\section{Definição do problema}

O problema proposto apresenta a trajetória de um projétil lançado por um obuseiro leve L117. Porém, ao invés de fornecer uma função explícita da altura em função do tempo, o enunciado fornece apenas alguns pontos conhecidos da trajetória. Os dados são os seguintes:

\begin{center}
\begin{tabular}{c|c}
Tempo (s) & Altura (m)\\
\hline
0 & 1,5\\
3 & 1007\\
7 & 2075\\
10 & 2670\\
14 & 3190\\
29 & 2339\\
31 & 1892\\
\end{tabular}
\end{center}

Como temos apenas valores discretos, precisamos construir uma função aproximada que represente a trajetória dentro do intervalo observado. Para isso, foi escolhido o método de interpolação de Newton, que constrói um polinômio que passa exatamente pelos pontos fornecidos.

A partir deste polinômio, o objetivo passa a ser encontrar o instante em que a altura é máxima. Isto pode ser feito procurando o ponto em que a derivada do polinômio é igual a zero.

\section{Método de interpolação}

A interpolação de Newton escreve o polinômio na forma:
\[
\begin{aligned}
P(x) ={}& b_0
+ b_1(x-x_0)
+ b_2(x-x_0)(x-x_1)\\
&+ \cdots
+ b_n(x-x_0)\cdots(x-x_{n-1}).
\end{aligned}
\]

\noindent
Os valores $b_0, b_1, \ldots, b_n$ são os coeficientes do polinômio. Eles são obtidos a partir da tabela de diferenças divididas. A primeira coluna da tabela contém os valores originais de altura.

As demais colunas são calculadas usando a seguinte relação:

\vspace{0.3cm}
\[
\begin{aligned}
f[x_i,\ldots,x_{i+j}]
={}&
\frac{
 f[x_{i+1},\ldots,x_{i+j}]
 -
 f[x_i,\ldots,x_{i+j-1}]
}{x_{i+j}-x_i}
\end{aligned}
\]

\vspace{0.2cm}

Dessa forma, cada coluna depende da coluna anterior, até que se obtenha toda a tabela. Os coeficientes usados no polinômio são os valores da primeira linha desta tabela.

Para os pontos fornecidos, foram encontrados aproximadamente os seguintes coeficientes:

\begin{center}
\begin{tabular}{c|r@{$\,\times\,$}l}
Coeficiente & \multicolumn{2}{c}{Valor}\\
\hline
$b_0$ & 1.50 & $10^{0}$\\
$b_1$ & 3.35 & $10^{2}$\\
$b_2$ & -9.74 & $10^{0}$\\
$b_3$ & -7.14 & $10^{-3}$\\
$b_4$ & 8.19 & $10^{-4}$\\
$b_5$ & -3.80 & $10^{-5}$\\
$b_6$ & 1.74 & $10^{-6}$\\
\end{tabular}
\end{center}


\begin{center}
{\small\textit{Nota: Os coeficientes foram arredondados para facilitar a leitura.}}
\end{center}

Com estes coeficientes, obtemos um polinômio de grau seis. Este grau aparece naturalmente porque foram usados sete pontos. O polinômio não é escrito expandido no relatório porque sua forma expandida não ajuda muito na compreensão do método. A forma de Newton já mostra melhor como os coeficientes foram obtidos e como o polinômio é avaliado.

\section{Implementação}

A implementação foi feita em Python. Primeiro, os pontos foram lidos a partir de um arquivo contendo pares de tempo e altura. Depois, foi construída a tabela de diferenças divididas.

A ideia principal do algoritmo é construir a tabela coluna por coluna. A primeira coluna recebe as alturas originais, enquanto as colunas seguintes são calculadas a partir das diferenças divididas da coluna anterior. A construção da tabela de diferenças divididas é apresentada no \autoref{alg:diferencas-divididas}.

\begin{figure*}[t]
\centering
\begin{minipage}{0.75\textwidth}
\begin{algorithm}[H]
\caption{Construção da tabela de diferenças divididas}
\label{alg:diferencas-divididas}
\begin{algorithmic}[]
\State $n \gets \text{tamanho}(pontos)$
\For{$i \gets 0$ até $n-1$}
    \State $tabela[i][0] \gets pontos[i].altura$
\EndFor
\For{$j \gets 1$ até $n-1$}
    \For{$i \gets 0$ até $n-j-1$}
        \State $num \gets tabela[i+1][j-1]-tabela[i][j-1]$
        \State $den \gets pontos[i+j].tempo-pontos[i].tempo$
        \State $tabela[i][j] \gets num/den$
    \EndFor
\EndFor
\State \Return $tabela$
\end{algorithmic}
\end{algorithm}
\end{minipage}
\end{figure*}

Após a criação da tabela, os coeficientes foram extraídos da primeira linha. Em seguida foi criada uma função para avaliar o polinômio em qualquer instante $t$.

Um teste importante foi verificar se o polinômio realmente passava pelos pontos originais. Ou seja, para cada tempo da tabela, avaliou-se o polinômio e comparou-se o valor obtido com a altura original. Esta verificação é simples, mas ajuda bastante, pois uma tabela de diferenças divididas incorreta ainda poderia produzir um número final, mas esse número não teria significado.

\section{Estimativa da altura máxima}

Depois de construir o polinômio de aproximação $P(t)$, resta estimar a altura máxima do projétil. Como $P(t)$ aproxima a função altura, passaremos a escrever $h(t)$ para simplificar a notação. Como queremos encontrar um máximo, procuramos um ponto em que:
\[
h'(t)=0.
\]

Foi utilizado o método de Newton para encontrar uma raiz da derivada da função altura.

O método de Newton aproxima da raiz em passos na forma:
\[
x_{n+1}=x_n-\frac{f(x_n)}{f'(x_n)}.
\]

Neste problema, a fórmula geral é aplicada sobre a variável de tempo. A função cujas raízes desejamos encontrar é a derivada da altura. Portanto,
\[
f(t)=\frac{dh}{dt}
\]

 e sua derivada é
\[
f'(t)=\frac{d^2h}{dt^2}.
\]

Substituindo essas expressões na fórmula de Newton, obtemos:
\[
t_{n+1}=t_n-\frac{\frac{dh}{dt}(t_n)}{\frac{d^2h}{dt^2}(t_n)}.
\]

Na implementação, as derivadas foram aproximadas numericamente. Com um chute inicial próximo da região onde a altura parecia ser máxima ($t=14s$), o método convergiu para:
\[
t \approx 18,652473 \text{ s}.
\]

Avaliando $h(t)$ nesse instante, obteve-se:
\[
h(18,652473) \approx 3402,236046 \text{ m}.
\]

Portanto, a altura máxima estimada foi de aproximadamente $3402,24$ metros, atingida cerca de $18,65$ segundos após o disparo.

\section{Uso do ChatGPT}

Para a segunda parte do trabalho, foi feita uma consulta ao ChatGPT. A entrada fornecida foi a mesma tabela de tempo e altura apresentada na Seção~1.

Foi solicitado que a ferramenta estimasse em qual instante o projétil atingiria sua altura máxima e qual seria essa altura. A resposta obtida também usou interpolação para modelar a trajetória e procurou o máximo da função interpoladora.

O resultado retornado foi compatível com o resultado da implementação:

\begin{center}
\begin{tabular}{c|c}
Grandeza & Valor estimado\\
\hline
Tempo da altura máxima & $18,652473$ s\\
Altura máxima & $3402,236046$ m\\
\end{tabular}
\end{center}

\section{Comparação dos resultados}

A implementação própria e a consulta ao ChatGPT chegaram ao mesmo resultado numérico. A tabela abaixo resume a comparação:

\begin{center}
\resizebox{\columnwidth}{!}{%
\begin{tabular}{l|c|c}
Grandeza & Implementação & ChatGPT\\
\hline
Tempo máx. & $18,652473$ s & $18,652473$ s\\
Altura máx. & $3402,236046$ m & $3402,236046$ m\\
\end{tabular}%
}
\end{center}

Como os valores coincidiram, considera-se que o resultado obtido é consistente com o método usado. Além disso, o fato de o polinômio reproduzir os pontos originais reforça que a interpolação foi implementada corretamente.

Mesmo assim, é importante observar que o resultado depende diretamente da qualidade dos dados e do modelo escolhido. O polinômio passa pelos pontos fornecidos, mas isso não significa necessariamente que ele represente perfeitamente a trajetória física real do projétil. Ele é uma aproximação construída a partir dos dados disponíveis.

\section{Conclusão}

Neste trabalho foi utilizada a interpolação de Newton para estimar a altura máxima de um projétil a partir de valores conhecidos de tempo e altura. A tabela de diferenças divididas permitiu obter os coeficientes do polinômio, e a altura máxima foi estimada procurando o ponto em que a derivada do polinômio é $0$.

O resultado encontrado foi uma altura máxima de aproximadamente $3402,24$ metros no instante $18,65$ segundos. O mesmo resultado foi obtido através da consulta ao ChatGPT, quando a ferramenta recebeu a tabela de dados e a solicitação de estimar a altura máxima e o tempo correspondente.


\end{document}
